\setcounter{chapter}{3}
\chapter{Lebesgue内測度}

\paragraph{本章の目標}
Lebesgue流の可測性
\[
\text{外測度}=\text{内測度}
\]
を定義するため，
Lebesgue内測度を導入する．
ただし，内測度は外測度や位相の言葉を用いて定義される副次的な量であり，
Carathéodoryが整理した抽象理論では外測度が主役となることを後に説明する．

% \section{動機づけ}

% 外測度 $\lambda^*(A)$ は，
% 集合 $A$ を外側から可算個の区間で覆うために必要な最小コストであった．
% これは外側から見た面積である．
% しかし，本当に面積を表しているなら，
% 集合を内側から近似したときにも同じ値が出てほしい．

% 第3章では，Lebesgue外測度が開被覆の下限としても記述できることを見た．
% したがって，外側近似で開集合が主役なら，
% 内側近似でコンパクト集合が主役になるのは自然である．
% 実際，$\RR^d$ で内側近似に最も適しているのはコンパクト集合である．
% コンパクト集合は「閉じていて有界な，逃げ場のない集合」であり，
% 実際の図形を有限の情報で近似する際に扱いやすい．
% そこで，$A$ に含まれるコンパクト集合の外測度をできるだけ大きくしたものを
% 内測度と呼ぶ．

\section{定義}

\begin{definition}[Lebesgue内測度]
集合 $A \subset \RR^d$ に対し，Lebesgue内測度 $\lambda_*(A)$ を
\[
\lambda_*(A)
:=
\sup\{\lambda^*(K) \mid K \subset A,\ K \text{ はコンパクト}\}
\]
で定める．
\end{definition}

\begin{remark}
この定義は，Lebesgue外測度や，$\RR^d$ のコンパクト集合という位相構造を用いて記述されている．
%したがって，内測度は外測度よりも幾何学的である一方，抽象的な集合論だけでは定義できない．
    第3章では，Lebesgue外測度が開被覆の下限として言い換えられることを見た．位相論において，開集合の対照的な概念は（閉集合ではなく）コンパクト集合であり，この観点で内測度は外測度の対照的な量である．
%    測度と位相は異なる概念だが，Euclid空間の場合には位相構造を使ってこの点は後にCarathéodory流へ移行する動機となる．
\end{remark}

% \begin{remark}[基本集合による内側近似がうまくいかない理由]
% Jordan内測度との類推から，
% \[
% \sup\{m_J(E)\mid E \subset A,\ E \text{ は基本集合}\}
% \]
% のように定めたくなるのは自然である．
% さらに，基本集合の可算和を許して内側から近似したくなるかもしれない．
% しかし，基本集合 $E$ が $A$ に含まれるなら，
% $E$ は開集合を含むので実際には
% \[
% E \subset \intr A
% \]
% である．
% したがって，この方法で見えているのは本質的には $A$ そのものではなく
% \[
% \intr A
% \]
% だけである．

% 特に，内部が空であっても大きさをもつ集合ではこの方法は破綻する．
% 例えば，内部が空で正の長さをもつfat Cantor集合では，
% 基本集合による内側近似は常に $0$ しか与えない．
% このため，Lebesgue内測度では基本集合ではなく，
% コンパクト集合による packing を用いる．
% \end{remark}

\begin{remark}[区間による内側近似がうまくいかない理由]
Jordan内測度との類推から，内測度を
\[
\gamma(A) := 
\sup\left\{ \sum_{i=1}^\infty |R_i| \;\middle|\; \bigsqcup_{i=1}^\infty R_i \subset A,\ R_i \text{ は互いに素な左半開区間}\right\}
\]
のように定めても，あまりうまくいかない．

まず，内側からの近似では，外測度のような（区間の重なりを許す） covering ではなく packing にとる必要がある．
区間の重なりを許すと同じ部分を何度でも数えられ，supremum が無限大になるからである．

しかし，互いに素な区間による packing にしても，左半開区間だけでは $A$ の内部しか見えない．
非退化な左半開区間 $R_i$ は空でない開集合を含むので，$R_i \subset A$ なら
\[
\intr R_i \subset \intr A
\]
である．
しかも境界は測度 $0$ だから，$|R_i|=|\intr R_i|$ であり，
この方法で得られる量は高々 $\lambda^*(\intr A)$ である．
したがって，$\intr A=\emptyset$ ならこの方法は常に $0$ しか与えない．

%伝統的に使われるLebesgue測度の正則性は，可測集合を内側から近似する基本対象がコンパクト集合であることを述べている．
そして，fat Cantor 集合（Smith--Volterra--Cantor 集合）のように，
内部が空でありながら正のLebesgue測度をもつコンパクト集合が存在するので，
区間 packing ではこのような集合を測れない．
このような背景を踏まえ，Lebesgue内測度では区間ではなく，コンパクト集合による packing を用いる．
\end{remark}

\begin{remark}
有界集合 $A$ をある有界閉区間 $I$ に入れて
\[
|I|-\lambda^*(I\setminus A)
\]
によって内測度を定義する流儀もある．
この定義が上のコンパクト集合による定義と一致することは，
開集合とコンパクト集合の可測性を示した後で証明する．
\end{remark}

\section{基本性質}

\begin{proposition}
任意の $A,B \subset \RR^d$ に対して次が成り立つ．
\begin{enumerate}
    \item $0 \le \lambda_*(A) \le \lambda^*(A)$
    \item $A \subset B$ ならば $\lambda_*(A) \le \lambda_*(B)$
\end{enumerate}
\end{proposition}
\begin{proof}
\begin{enumerate}
    \item コンパクト集合 $K \subset A$ に対して，外測度の単調性より
    \[
    \lambda^*(K) \le \lambda^*(A)
    \]
    である．
    左辺について上限を取れば
    \[
    \lambda_*(A) \le \lambda^*(A)
    \]
    を得る．
    また外測度は常に非負なので，その上限である $\lambda_*(A)$ も非負である．

    \item $A \subset B$ とする．
    $K \subset A$ なる任意のコンパクト集合は，そのまま $K \subset B$ も満たす．
    したがって，$A$ に対して取る上限は $B$ に対して取る上限以下である．
\end{enumerate}
\end{proof}

\begin{proposition}[コンパクト集合では内外が一致する]
コンパクト集合 $K \subset \RR^d$ に対して
\[
\lambda_*(K)=\lambda^*(K)
\]
が成り立つ．
\end{proposition}
\begin{proof}
$K$ 自身がコンパクトだから，
内測度の上限を取る集合の中に $K$ そのものが含まれている．
したがって
\[
\lambda_*(K) \ge \lambda^*(K)
\]
である．
逆向きの不等式は前命題から従う．
\end{proof}

\begin{corollary}[外測度零集合の内測度]
$\lambda^*(A)=0$ ならば
\[
\lambda_*(A)=0
\]
である．
\end{corollary}
\begin{proof}
基本性質の (1) より
\[
0 \le \lambda_*(A) \le \lambda^*(A)=0
\]
だからである．
\end{proof}

\section{具体例}

\begin{example}
\begin{enumerate}
    \item 任意のコンパクト集合では内測度と外測度が一致する
    \item 任意の可算集合 $A$ に対して
    \[
    \lambda_*(A)=0
    \]
    である
    \item 特に
    \[
    \QQ \cap [0,1]
    \]
    は内測度 $0$ をもつ
\end{enumerate}
\end{example}

\begin{proof}
\begin{enumerate}
    \item 直前の命題で示した．
    \item 第3章で可算集合の外測度は $0$ であることを示したから，
    直前の系より
    \[
    \lambda_*(A)=0
    \]
    である．
    \item (2) を
    \[
    A=\QQ \cap [0,1]
    \]
    に適用すればよい．
\end{enumerate}
\end{proof}

\section{解釈}

Lebesgue内測度 $\lambda_*(A)$ は，
「$A$ の中に入るコンパクト集合のうち，どこまで大きな面積を確保できるか」
を表している．
したがって
\[
\lambda^*(A)=\lambda_*(A)
\]
とは，
外から見た大きさと内から見た大きさが一致して，
その集合の面積がぶれないことを意味している．
この等式を次章でLebesgue可測性の定義として採用する．

ただし，この定義はまだ補助的なものである．
理由は2つある．
\begin{enumerate}
    \item コンパクト集合という位相的概念に依存している
    \item この定義だけから可測集合全体の構造を調べるのは難しい
\end{enumerate}
第5章では，Lebesgue測度が実際にどのような性質をもつか，
特に正則性
\[
\lambda(A)
=
\sup\{\lambda(K)\mid K \subset A,\ K \text{ compact}\}
\]
がどのように内測度と結びつくかを見る．
そして第6章では，内測度を使わずに測度を構成できる
Carathéodory流の定義へ進む．

